\section{Limitations}
\label{sec:limitations}

We state the boundaries of the evidence plainly, because a negative result is only as
useful as its scope is honest.

\textbf{Scope of the negatives.} The three studies are scoped, not universal. Q17 covers
two assets and eight paired dates under a quote-based utility that uses quote endpoints and
assumed costs. It does not observe realized fills, endogenous market impact, funding, or
queue priority, so it is a decision-utility statement, not a live-trading profitability
statement. Q18 uses sampled-L2 receipt-clock data over nine dates. Its date panel is small
and its intervals are wide, so it rejects a simple sufficiency rule rather than proving
that depth never helps. Q16's reconstruction diagnostic uses development-exposed dates and
is descriptive rather than a pristine prospective holdout.

\textbf{Data realism.} Provider receipt time is uncalibrated against exchange time, and the
sampled-L2 feed does not establish native queue priority or calibrated latency. Where an
entity or feed could not be certified, the corresponding claim is withheld. The
native-FIFO Q16 endpoint (Section~\ref{sec:synthesis}) is the clearest such withholding.

\textbf{Estimation confounds.} Several neural fits reached their fixed iteration budget,
and their warnings are preserved rather than suppressed. Full-depth neural underperformance
can therefore reflect an optimization limit rather than an information limit. Because
changing only that family cannot rescue a result that requires agreement across families,
we do not lean on any single neural cell.

\textbf{What we do not claim.} We do not claim that information is worthless, that richer
feeds should never be purchased, or that these negatives generalize to other venues,
horizons, or asset classes. We claim that, measured at the decision with everything else
held fixed, the value of more and better-reconstructed market information was small,
conditional, and fragile to the policy and cost layer, across three questions on real
data.
